\documentclass[UTF8,11pt]{ctexart}
\usepackage[margin=1.02in]{geometry}
\usepackage{amsmath,amssymb,amsthm,mathtools}
\usepackage{graphicx}
\usepackage{booktabs}
\usepackage{enumitem}
\usepackage[round,authoryear]{natbib}
\usepackage[hidelinks,hypertexnames=false]{hyperref}
\usepackage[nameinlink,capitalize,noabbrev]{cleveref}
\usepackage{microtype}
\usepackage{xcolor}
\usepackage{float}
\emergencystretch=2em

\newtheorem{theorem}{定理}[section]
\newtheorem{proposition}[theorem]{命题}
\newtheorem{corollary}[theorem]{推论}
\newtheorem{lemma}[theorem]{引理}
\theoremstyle{definition}
\newtheorem{definition}[theorem]{定义}
\newtheorem{remark}[theorem]{注}

\crefname{theorem}{定理}{定理}
\crefname{proposition}{命题}{命题}
\crefname{corollary}{推论}{推论}
\crefname{lemma}{引理}{引理}
\crefname{definition}{定义}{定义}
\crefname{remark}{注}{注}

\newcommand{\R}{\mathbb R}
\newcommand{\Z}{\mathbb Z}
\newcommand{\N}{\mathbb N}
\newcommand{\E}{\mathbb E}
\newcommand{\diff}{\mathop{}\!d}

\title{补偿性当前选择与反事实识别}
\author{}
\date{2026年8月}

\begin{document}
\maketitle

\begin{abstract}
对一个有限的当前变化进行精确补偿转移，可以把该变化置于共同的货币尺度上；令变化幅度趋近于零，则可揭示相应的局部斜率要素。当目标环境本身可以在局部被观察时，在不同继承状态下重复这一比较即可得到直接的比较静态。本文关注的是反事实情形：目标环境中的基准选择可以观察，但直接计算比较静态所需的邻近补偿比较并不存在。

观测到的比较构成一个有向图。其当前层投影中的循环会消去任意的当前节点价值。在消去条件成立时，转移量具有加法表示，而数据所识别的延续差异恰好是当前循环空间在延续关联映射下的像。其维数为
\[
\operatorname{rank}\!\begin{pmatrix}D\\Q\end{pmatrix}-\operatorname{rank}D.
\]
新增一条比较至多增加一个延续方向，并且当且仅当它在当前关联上是冗余的、而在当前--延续联合关联上并不冗余时才会增加一个方向。随后，一个原始的双区块设计可以识别 $H_{xa}$ 与 $H_{xx}$，尽管任一抽样点上的水平值 $H(z)$ 都没有得到点识别。更一般地，收缩循环可以消去所有低阶多项式项而保留某个高阶矩，因此可以在不恢复低阶 jet 的情况下识别高阶局部泛函。这些已识别导数进入目标环境的一阶条件，从而识别反事实比较静态。对多个环境以及状态动态的扩展，则在基于比较的识别步骤之后使用标准的势函数和动力系统论证。
\end{abstract}

\section{引言}

一个当前选择既可能改变眼下的结果，也可能改变其后续后果。在行动 $x$ 和继承状态 $a$ 下，将当前方案与另一个方案比较：后者包含额外的 $\Delta$ 单位行动，但少了 $p\Delta$ 单位的货币计价物。记使两者无差异的补偿率为 $p_\Delta(x,a)$。在共同货币尺度上，
\[
p_\Delta(x,a)
=
\frac{W(x+\Delta,a)-W(x,a)}{\Delta},
\]
其中 $W$ 包含当前行动所引致的延续后果。因此极限补偿率
\[
\pi(x,a):=\lim_{\Delta\downarrow0}p_\Delta(x,a)=W_x(x,a)
\]
就是用计价物表示的 Samuelson 局部斜率要素。图~\ref{fig:primitive} 给出了这一比较及其极限。

令 $C(x,a)$ 为已知的当前资源成本。在一个内部严格局部最优点，
\[
\pi(x,a)=C_x(x,a).
\]
在不同继承状态下重复这些比较，就可以描出一条正则的局部选择分支 $x=g(a)$，从而得到直接比较静态。图~\ref{fig:locus} 给出了这一基准情形。

一个两阶段保险合同提供了方便的贯穿例子。当前后果记录当期保障及其他合同条款，其中所有已知的货币金额都已先行净除；延续后果则是一份状态依赖的续保合同。一个诱导菜单可以保持当前后果不变，同时改变续保条款，并通过当期保费转移抵消差异。因此，即使某一保险公司的目标菜单把当前保障与唯一的续保规则绑定，同一个当前节点在实验域中仍可能对应多个延续标签。下文理论研究的就是这类已观察复合比较中的信息；它并不要求目标菜单本身包含所有可能的重组。

现在令环境 $A$ 在局部被观察，而环境 $B$ 为目标环境：我们观察到其基准选择，但没有邻近补偿比较。记
\[
H(x,a):=(V\circ\Gamma_A)(x,a)-(V\circ\Gamma_B)(x,a).
\]
目标响应取决于 $H$ 的导数。物理增量之和为零，并不能消去一个不受限制的当前价值。精确消去要求每一个当前节点都达到平衡。若 $D$ 为节点--边关联矩阵，则满足
\[
Dr=0
\]
的边向量 $r$ 是一个循环；等价地，它属于图的循环空间。此时当前节点价值沿环路望远镜相消，但延续载荷 $Qr$ 不必消失。图~\ref{fig:cancellation} 展示了闭合的当前投影与不闭合的延续提升。

在一个可观察的消去条件下，转移量具有加法表示。消去当前部分后仍然存活的延续信息为
\[
Q(\ker D).
\]
其维数等于把延续关联加入当前关联之后新增的秩。新增一条比较恰好在如下情况下增加一个延续方向：它的当前关联已经被已有比较张成，但其当前--延续联合关联并未被已有比较张成。因此，比较支持所包含的信息可以逐边归属。

这种有限几何有一个比恢复水平值更强的局部后果。由两个四边循环构成的原始设计可以识别 $H_{xa}(z_0)$ 与 $H_{xx}(z_0)$，同时每一个抽样值 $H(z)$ 仍然未识别。更一般地，收缩循环权重可以消去所有低于 $m$ 阶的多项式项，却保留一个非零的 $m$ 阶矩。于是相应的 $m$ 阶微分泛函能够被识别，而任何非零的低阶 jet 泛函仍未识别。反事实一阶条件把这两个二阶延续修正转换为目标比较静态。

本文论证顺序遵循显示偏好的传统：先有观测比较和有限一致性限制，再有数值表示。\citet{samuelson1938,samuelson1948} 以这种方式构造消费者理论，\citet{samuelson1950} 则把局部显示的价格比称为偏好场的斜率要素。\citet{debreu1954} 区分偏好排序与其数值表示，\citet{afriat1967} 则从有限支出数据构造效用。在部分可观测情形下，\citet{gorno2019} 研究偏好的识别，\citet{chambersetal2021} 研究从不断扩展的有限二元实验中恢复偏好。本文所识别的对象更窄：足以用于一个反事实比较静态的延续泛函。

货币度量效用与补偿需求由 \citet{fuchsseliger1989,fuchsseliger1990} 研究；\citet{bhattacharya2015} 给出了离散选择下非参数的货币度量福利结果。准线性可理性化性与可积性可以用循环限制来刻画，参见 \citet{rochet1987,browncalsamiglia2007,nockeschutz2017}；\citet{castillofreer2023} 使用有向图构造准线性理性化。\citet{makowskiostroy2013} 用线性空间对偶和循环单调性表述偏好揭示。下文的图关系具有不同作用：循环用于消去不受限制的当前节点价值并识别延续差异。加法表示则属于更一般的消去条件与加法测量理论 \citep{kraftprattseidenberg1959,scott1964,doignonfalmagne1974,fishburnetal1988,fishburn2001}。

本文维持的货币结构是可转移的计价物，以及每个观测二元比较都存在唯一的精确补偿转移。主要识别结果之后的扩展使用标准的图势函数、迭代映射术语和递归方法 \citep{elaydi2005,kuznetsov2004,stokeylucasprescott1989}。

\section{原始补偿二元实验}
\label{sec:primitive}

令 $x$ 表示标量当前行动，$a$ 表示继承状态，$m$ 表示货币计价物。对于每一个 $(x,a)$，令 $W(x,a)$ 表示在状态 $a$ 下选择 $x$ 所产生的非计价物后果的货币度量价值。因此当前方案 $(x,m)$ 的价值为
\begin{equation}
 m+W(x,a).
 \label{eq:total-value}
\end{equation}
当前行动引致的所有延续后果均包含在 $W$ 中。

固定 $(x,a)$ 和 $\Delta>0$。比较
\begin{equation}
 B_0(x)=(x,m)
 \label{eq:B0}
\end{equation}
与
\begin{equation}
 B_1(x,a;p,\Delta)=(x+\Delta,m-p\Delta),
 \label{eq:B1}
\end{equation}
其中后者的延续后果由 $(x+\Delta,a)$ 所引致。其他所有当前属性保持不变，只改变 $p$。

更一般地，把一个复合方案写成
\[
(\xi,y,m),
\]
其中 $\xi$ 是当前后果，$y$ 是延续后果。对观察者比较域中的任意一对方案，通过只改变其中一端的货币量来定义精确补偿转移。在保险例子中，$y$ 是续保条款，货币调整则是保费转移。上面的邻近行动比较是当前行动同时决定其延续后果的特殊情形。

\begin{figure}[H]
\centering
\refstepcounter{figure}\label{fig:primitive}
\includegraphics[width=\textwidth]{fig1_colored.jpg}
\par\smallskip\small\textit{图示说明：左图通过改变补偿率 $p$ 得到有限切换率 $p_\Delta(x,a)$；右图给出 $(x,m)$ 空间中的二元方案几何。令 $\Delta\downarrow0$ 后，割线斜率趋于局部斜率 $-\pi(x,a)$。}
\end{figure}

对每个 $\Delta$，令 $p_\Delta(x,a)$ 由下式定义：
\begin{equation}
B_1(x,a;p_\Delta(x,a),\Delta)\sim B_0(x),
\label{eq:threshold-indifference}
\end{equation}
并令
\[
\tau_\Delta(x,a):=p_\Delta(x,a)\Delta
\]
为精确补偿转移。按照图~\ref{fig:primitive} 的方向约定，
\begin{align}
p<p_\Delta(x,a)&\quad\Longrightarrow\quad B_1(x,a;p,\Delta)\succ B_0(x),
\label{eq:below-threshold}\\
p=p_\Delta(x,a)&\quad\Longrightarrow\quad B_1(x,a;p,\Delta)\sim B_0(x),
\label{eq:at-threshold}\\
p>p_\Delta(x,a)&\quad\Longrightarrow\quad B_1(x,a;p,\Delta)\prec B_0(x).
\label{eq:above-threshold}
\end{align}
在补偿率处，
\begin{equation}
 m+W(x,a)
 =m-p_\Delta(x,a)\Delta+W(x+\Delta,a),
\end{equation}
从而
\begin{equation}
\boxed{
 p_\Delta(x,a)
 =
 \frac{W(x+\Delta,a)-W(x,a)}{\Delta}.}
\label{eq:finite-rate-secant}
\end{equation}
因此 $p_\Delta(x,a)$ 是单位有限当前增量对应的补偿变差。

沿用 Samuelson 的“斜率要素”术语，定义极限补偿率
\begin{equation}
\boxed{
\pi(x,a):=\lim_{\Delta\downarrow0}p_\Delta(x,a).}
\label{eq:pi-def}
\end{equation}
若 $W(\cdot,a)$ 在 $x$ 处可微，则式~\eqref{eq:finite-rate-secant} 是其右差商，因此
\begin{equation}
\boxed{W_x(x,a)=\pi(x,a).}
\label{eq:Wx-pi}
\end{equation}
若 $\pi$ 在 $(x,a)$ 附近连续可微，则
\begin{equation}
W_{xa}=\pi_a,
\qquad
W_{xx}=\pi_x.
\label{eq:pi-derivatives}
\end{equation}

\section{局部选择函数}
\label{sec:locus}

令当前选择问题的净货币价值为 $W(x,a)-C(x,a)$，其中 $C(x,a)$ 是继承状态 $a$ 下行动 $x$ 的资源成本。内部最优解满足
\begin{equation}
W_x(x,a)=C_x(x,a).
\label{eq:FOC-WC}
\end{equation}
由式~\eqref{eq:Wx-pi}，在能够进行补偿实验的地方，这个条件是可观察的：
\begin{equation}
\boxed{
\pi(x,a)=C_x(x,a).}
\label{eq:FOC-pi}
\end{equation}
定义
\begin{equation}
F(x,a):=\pi(x,a)-C_x(x,a).
\label{eq:F-def}
\end{equation}
$F$ 的正则零水平集定义局部选择函数。

\begin{figure}[H]
\centering
\refstepcounter{figure}\label{fig:locus}
\includegraphics[width=\textwidth]{fig2_locus.jpg}
\par\smallskip\small\textit{图示说明：在每个继承状态 $a$ 下重复图~\ref{fig:primitive} 的补偿二元实验，满足 $\pi(x,a)=C_x(x,a)$ 的点构成最优当前选择轨迹 $x=g(a)$；其局部斜率 $g'(a_0)$ 即直接比较静态。}
\end{figure}

假设 $F$ 在 $(x_0,a_0)$ 附近连续可微，$F(x_0,a_0)=0$，并且
\begin{equation}
C_{xx}(x_0,a_0)-\pi_x(x_0,a_0)>0.
\label{eq:regularity-F}
\end{equation}
隐函数定理给出唯一的可微局部选择分支 $x=g(a)$，满足
\begin{equation}
\boxed{F(g(a),a)=0}
\label{eq:g-locus}
\end{equation}
以及
\begin{equation}
\boxed{
 g'(a)
 =
 \frac{\pi_a(g(a),a)-C_{xa}(g(a),a)}
 {C_{xx}(g(a),a)-\pi_x(g(a),a)}.}
\label{eq:direct-gprime}
\end{equation}
当邻近补偿比较可观察时，式~\eqref{eq:direct-gprime} 就是直接比较静态。

\subsection{反事实问题}

设环境 $A$ 中可以观察邻近补偿比较；在环境 $B$ 中，只观察到基准当前选择 $g_B(a_0)$，但观察不到邻近补偿比较。写为
\[
W_A=u_A+V\circ\Gamma_A,
\qquad
W_B=u_B+V\circ\Gamma_B.
\]
这一记号定义了本文关心的结构性反事实对象。这里并不使用 $V$ 的任何数值来理性化有限比较数据。相反，第~\ref{sec:cancellation} 节刻画观测到的补偿转移何时自身具有当前--延续加法表示。当结构模型被维持时，其延续价值在观测标签上给出这种表示的一个实例。

环境 $A$ 中的补偿观测识别 $W_A$ 的局部导数。目标响应还取决于
\[
H(x,a):=(V\circ\Gamma_A)(x,a)-(V\circ\Gamma_B)(x,a)
\]
的导数。问题是：如何从已观测的补偿比较中识别 $H$ 的有限与局部泛函。

\section{比较图与消去条件}
\label{sec:cancellation}

考虑有限个已观测补偿二元比较组成的集合 $E$。按照实验协议，所有已知且可用货币计量的当前侧金额都先从原始补偿转移中净除。记调整后的精确转移为 $\tau_e$。令
\[
\mathcal X=\{\xi_0,\ldots,\xi_N\}
\]
为完成这一已知调整后在这些比较中出现的不同\emph{当前节点}。一个当前节点记录当前侧仍具有不受限制货币价值的属性。因此，即使某个已知的环境特定当前项不同，只要两个方案具有相同当前组合，它们仍可以使用同一个节点；该已知项已经从 $\tau_e$ 中净除。每个比较 $e$ 有尾节点 $s(e)$ 与头节点 $t(e)$。

令 $D\in\Z^{(N+1)\times |E|}$ 为有向节点--边关联矩阵，其第 $e$ 列为
\begin{equation}
 d_e=\mathbf e_{t(e)}-\mathbf e_{s(e)}.
 \label{eq:current-incidence}
\end{equation}
因此 $Dz=0$ 表示整数边集合 $z\in\Z^{E}$ 在每个当前节点上的流入次数与流出次数完全相同；也就是说，$Dz=0$ 是精确节点平衡。

实验协议还记录每条比较所携带的带符号延续后果。固定有限个延续标签并做一个归一化后，把边 $e$ 的延续对比写成 $q_e\in\Z^m$，并把这些列收集为
\[
Q\in\Z^{m\times |E|}.
\]
令 $\tau=(\tau_e)_{e\in E}\in\R^{E}$。

\begin{definition}[消去条件]
\label{def:joint-cancellation}
若
\begin{equation}
\boxed{
Dz=0,\quad Qz=0
\quad\Longrightarrow\quad
z^{\mathsf T}\tau=0
\qquad(z\in\Z^{E}),}
\label{eq:joint-cancellation}
\end{equation}
则称有限补偿数据满足\emph{消去条件}。
\end{definition}

该条件要求：只要当前关联与延续对比同时为零，货币转移的加总也必须为零。

\begin{theorem}[加法货币表示]
\label{thm:additive-graph}
消去条件成立，当且仅当存在当前节点势 $u\in\R^{N+1}$ 和延续势 $v\in\R^m$，使得
\begin{equation}
\boxed{
\tau=D^{\mathsf T}u+Q^{\mathsf T}v.}
\label{eq:graph-representation}
\end{equation}
该表示在观测关联矩阵与延续矩阵的零空间意义下唯一。势 $(u,v)$ 是由有限数据得到的表示坐标，并不是第~\ref{sec:locus} 节引入的结构函数 $(u_J,V)$。
\end{theorem}

图~\ref{fig:cancellation} 展示一个节点关联循环及其延续提升。

\begin{figure}[H]
\centering
\refstepcounter{figure}\label{fig:cancellation}
\includegraphics[width=\textwidth]{fig3_cancellation.jpg}
\par\smallskip\small\textit{图示说明：面板 (i) 给出满足 $Dr=0$ 的当前层关联循环；面板 (ii) 把同一组基点提升到延续空间，提升后的路径一般不闭合，其垂直缺口即可观察的延续差异；面板 (iii) 令循环尺度收缩到零，从而把有限循环转化为局部导数信息并最终进入目标一阶条件。}
\end{figure}

\begin{corollary}[循环方程]
\label{cor:cycle-residual}
在消去条件成立时，每个满足 $Dz=0$ 的当前循环 $z\in\Z^E$ 都满足
\begin{equation}
\boxed{
z^{\mathsf T}\tau=(Qz)^{\mathsf T}v.}
\label{eq:cycle-residual}
\end{equation}
\end{corollary}

由于 $z^{\mathsf T}D^{\mathsf T}u=(Dz)^{\mathsf T}u=0$，当前节点价值会从每一条循环方程中消失。

\begin{theorem}[有限识别]
\label{thm:feasible-finite}
假设消去条件成立，并考虑模型集合
\[
\mathcal M(\tau)
=
\{(u,v):D^{\mathsf T}u+Q^{\mathsf T}v=\tau\}.
\]
对目标延续对比 $c\in\Z^m$，$c^{\mathsf T}v$ 在 $\mathcal M(\tau)$ 上得到点识别，当且仅当
\begin{equation}
\boxed{
\exists\,k\ge1,\ z\in\Z^E
\quad\text{使得}\quad
Dz=0,\qquad Qz=kc.}
\label{eq:feasible-cycle}
\end{equation}
当式~\eqref{eq:feasible-cycle} 成立时，
\begin{equation}
\boxed{
c^{\mathsf T}v=\frac{z^{\mathsf T}\tau}{k}.}
\label{eq:feasible-price}
\end{equation}
否则，相对于式~\eqref{eq:graph-representation}，
\begin{equation}
\boxed{
\{c^{\mathsf T}v:(u,v)\in\mathcal M(\tau)\}=\R.}
\label{eq:feasible-allR}
\end{equation}
\end{theorem}

被识别的延续对比，恰好是整数循环格在有限复制之后经过 $Q$ 映射得到的像。

除非另有说明，下文的有限识别均相对于加法表示~\eqref{eq:graph-representation} 而言；局部识别则相对于满足观测循环方程的约化式 $C^m$ 延续函数。若固定的共同 $V$ 与固定延续映射在不同点之间施加进一步结构限制，则相应的结构识别集可能更小。

\section{由收缩循环得到的局部识别}
\label{sec:shrinking}

令 $z=(x,a)$，并固定 $z_0=(x_0,a_0)$。在局部设计 $n$ 中，令
\[
z_{1n},\ldots,z_{M_n n}
\]
为抽样设计点，并写
\[
A(z):=(V\circ\Gamma_A)(z),
\qquad
B(z):=(V\circ\Gamma_B)(z),
\qquad
H(z):=A(z)-B(z).
\]
把 $Q_n$ 中的延续标签按如下次序排列：
\[
A(z_{1n}),\ldots,A(z_{M_n n}),
B(z_{1n}),\ldots,B(z_{M_n n}).
\]
当有限表示通过结构模型解释时，延续势 $v$（在共同归一化下）把这些标签分别赋值为 $A(z_{in})$ 与 $B(z_{in})$。为了给出锐利的局部结论，我们直接使用满足观测循环方程的约化式函数 $A$ 与 $B$。设计 $n$ 中的原始补偿比较具有观测当前关联 $D_n$、延续对比 $Q_n$ 与精确转移 $\tau_n$。

\subsection{目标环境与观测比较}
\label{sec:environment-support}

目标环境在 $z_0$ 的一个邻域上定义为
\[
W_B(x,a)=u_B(x,a)+V(\Gamma_B(x,a)).
\]
令 $\mathcal E_n$ 表示局部设计 $n$ 中观测到的补偿比较。族 $\{\mathcal E_n\}$ 不必包含由环境 $B$ 中所有邻近方案所组成的比较。它还可以包含这样的诱导菜单：把一个当前后果 $\xi$ 与某个延续合同 $\Gamma_J(z)$ 配对，即使该复合方案并不属于目标选择集。在保险例子中，这类合同变体保持当前保障不变，却替换成不同的状态依赖续保条款。下文的识别均相对于 $\{\mathcal E_n\}$ 而言；$\Gamma_B$ 仍在完整的局部选择集上有定义。

\subsection{来自循环空间的延续对比}

令
\[
S_n:=(I_{M_n}\ I_{M_n}),
\qquad
T_n:=(I_{M_n}\ 0).
\]
与 $A$--$B$ 对比相关的整数循环格为
\begin{equation}
\mathcal C_n^{\mathbb Z}
:=
\left\{
r\in\mathbb Z^{E_n}:
D_nr=0,\ S_nQ_nr=0
\right\}.
\label{eq:integer-local-cycles}
\end{equation}
它们刻画观测比较之间的有限复制与消去关系。

这些观测方程可以线性组合。定义实循环空间
\begin{equation}
\mathcal C_n
:=
\operatorname{span}_{\mathbb R}\mathcal C_n^{\mathbb Z}
=
\left\{
r\in\mathbb R^{E_n}:
D_nr=0,\ S_nQ_nr=0
\right\}.
\label{eq:real-local-cycles}
\end{equation}
之所以有上述等式，是因为 $D_n$ 与 $Q_n$ 的系数都是整数：有理零空间向量张成实零空间，而通过清除分母可把它们转化为整数循环。

对 $r\in\mathcal C_n$，总延续对比具有形式
\[
Q_nr=(\lambda,-\lambda),
\qquad
\lambda=T_nQ_nr.
\]
由此得到的对比空间为
\begin{equation}
\boxed{
\Lambda_n:=T_nQ_n\mathcal C_n
\subseteq\mathbb R^{M_n}.}
\label{eq:gap-weight-space}
\end{equation}
在消去条件下，每个 $r\in\mathcal C_n^{\mathbb Z}$ 都满足
\[
r^{\mathsf T}\tau_n
=
\sum_i\lambda_iH(z_{in}).
\]
由线性性，同一恒等式对任意 $r\in\mathcal C_n$ 也成立：
\begin{equation}
r^{\mathsf T}\tau_n
=
\sum_i\lambda_iH(z_{in}),
\qquad
Q_nr=(\lambda,-\lambda).
\label{eq:gap-cycle-price}
\end{equation}

\subsection{收缩循环与有限差分}
\label{sec:finite-differences}

固定包含局部设计点的 $z_0$ 开邻域 $U$。若一个光滑扰动 $q$ 满足
\begin{equation}
\sum_i\lambda_iq(z_{in})=0
\qquad
\text{对所有 }n\text{ 以及所有 }\lambda\in\Lambda_n,
\label{eq:null-perturbation}
\end{equation}
则称其在观测上为零。约化式光滑函数类上的识别具有通常的零空间形式：一个局部线性泛函得到点识别，当且仅当它在所有观测为零的扰动上都为零。下文只用这一事实来说明锐利性；具体识别公式则直接来自有限差分。

对 $t=(t_0,t_x,t_a,t_{xx},t_{xa},t_{aa})\in\R^6$，定义
\begin{equation}
\mathcal L_t f
:=
t_0f+t_xf_x+t_af_a+t_{xx}f_{xx}+t_{xa}f_{xa}+t_{aa}f_{aa}
\quad\text{在 }z_0\text{ 处}.
\label{eq:Lt-definition}
\end{equation}
令 $P_n$ 的第 $i$ 行为
\[
\left(
1,\delta x_{in},\delta a_{in},
\frac12\delta x_{in}^2,
\delta x_{in}\delta a_{in},
\frac12\delta a_{in}^2
\right),
\]
其中 $(\delta x_{in},\delta a_{in})=z_{in}-z_0$。

\begin{theorem}[收缩有限差分]
\label{thm:primitive-local}
假设支持半径趋于零。令 $\lambda_n\in\Lambda_n$ 与 $\alpha_n\ne0$ 满足
\begin{equation}
\frac{P_n^{\mathsf T}\lambda_n}{\alpha_n}\longrightarrow t
\label{eq:moment-limit}
\end{equation}
以及
\begin{equation}
\frac1{|\alpha_n|}\sum_i|\lambda_{in}|
(\delta x_{in}^2+\delta a_{in}^2)=O(1).
\label{eq:moment-variation}
\end{equation}
那么，对 $z_0$ 邻域中的任意 $C^2$ 函数 $f$，
\begin{equation}
\boxed{
\frac1{\alpha_n}\sum_i\lambda_{in}f(z_{in})
\longrightarrow \mathcal L_tf(z_0).}
\label{eq:cycle-moment-limit}
\end{equation}
在消去条件下，当 $f=H$ 时，左端可由补偿转移的标准化循环和直接观察。因此 $\mathcal L_tH(z_0)$ 得到点识别。
\end{theorem}

\begin{proof}
Taylor 展开给出
\[
f(z_{in})
=P_{n,i\cdot}\,j^2_{z_0}f+R_{in},
\qquad
R_{in}=o(\delta x_{in}^2+\delta a_{in}^2),
\]
并且该余项在收缩支持上一致成立。式~\eqref{eq:moment-limit}--\eqref{eq:moment-variation} 给出极限；式~\eqref{eq:gap-cycle-price} 则使其在 $H$ 上可观察。
\end{proof}

熟悉的中心二阶差分与混合矩形差分分别对应
\[
(1,-2,1)/h^2
\quad\text{和}\quad
(1,-1,-1,1)/(hk),
\]
前提是相应权重属于 $\Lambda_n$。

\subsection{不恢复水平值的二阶导数识别}
\label{sec:derivatives-without-levels}

比较图是观测复合方案在当前层的投影。令 $\mathcal A_n$ 表示设计 $n$ 中观测到的复合方案集合，并把典型元素写成 $(\xi,y)$，省略货币。每条观测比较边 $e$ 具有形式
\[
(\xi_{s(e)},y_e^-)
\quad\text{对比}\quad
(\xi_{t(e)},y_e^+),
\]
因此其当前关联为
\[
d_e=\mathbf e_{t(e)}-\mathbf e_{s(e)},
\]
延续对比为
\[
q_e=[y_e^+]-[y_e^-].
\]

\begin{proposition}[延续像空间的维数]
\label{prop:support-rank}
令 $D$ 与 $Q$ 分别为任意有限观测复合菜单支持的当前关联矩阵与延续关联矩阵，并定义
\[
\mathcal T:=Q(\ker D).
\]
则
\begin{equation}
\boxed{
\dim \mathcal T
=
\operatorname{rank}
\begin{pmatrix}
D\\ Q
\end{pmatrix}
-
\operatorname{rank}D.}
\label{eq:support-rank}
\end{equation}
因此，消去当前部分之后，观测支持携带非零延续信息，当且仅当
\[
\operatorname{rank}
\begin{pmatrix}
D\\ Q
\end{pmatrix}
>
\operatorname{rank}D.
\]

对第~\ref{sec:shrinking} 节中的局部 $A$--$B$ 对比，令
\[
M_n:=
\begin{pmatrix}
D_n\\ S_nQ_n
\end{pmatrix}.
\]
则
\begin{equation}
\boxed{
\dim\Lambda_n
=
\operatorname{rank}
\begin{pmatrix}
M_n\\ T_nQ_n
\end{pmatrix}
-
\operatorname{rank}M_n.}
\label{eq:gap-support-rank}
\end{equation}
因此，比较支持给出非零 $A$--$B$ 对比，恰好等价于：在已消除当前关联和共同延续关联之后，加入对比关联 $T_nQ_n$ 会提高秩。
\end{proposition}

\begin{proof}
$Q$ 在 $\ker D$ 上的限制，其核为 $\ker D\cap\ker Q$。秩--零化度定理给出
\[
\dim Q(\ker D)
=
\dim\ker D-
\dim(\ker D\cap\ker Q).
\]
由于
\[
\ker D\cap\ker Q
=
\ker
\begin{pmatrix}
D\\ Q
\end{pmatrix},
\]
式~\eqref{eq:support-rank} 成立。对 $T_nQ_n$ 在 $\ker M_n$ 上的限制应用同样论证，即得式~\eqref{eq:gap-support-rank}。
\end{proof}

\begin{theorem}[新增一条比较带来的秩增量]
\label{thm:marginal-edge}
令
\[
\nu(D,Q)
:=
\operatorname{rank}
\begin{pmatrix}
D\\Q
\end{pmatrix}
-
\operatorname{rank}D
=
\dim Q(\ker D).
\]
新增一条观测比较，其当前关联列为 $d$，延续关联列为 $q$。若这条比较引入新的当前节点或延续标签，则通过补零把旧矩阵嵌入扩大的坐标空间。记
\[
D^+=(D\ d),
\qquad
Q^+=(Q\ q).
\]
则
\begin{equation}
\boxed{
\nu(D^+,Q^+)-\nu(D,Q)\in\{0,1\}.}
\label{eq:marginal-info-01}
\end{equation}
该增量等于一，当且仅当
\begin{equation}
\boxed{
d\in\operatorname{col}D
\quad\text{且}\quad
\begin{pmatrix}d\\q\end{pmatrix}
\notin
\operatorname{col}
\begin{pmatrix}D\\Q\end{pmatrix}.}
\label{eq:marginal-info-condition}
\end{equation}
因此，一条比较恰好在“当前关联冗余、但当前--延续联合关联不冗余”时增加延续信息。

对局部 $A$--$B$ 对比空间，令
\[
m=
\begin{pmatrix}d\\Sq\end{pmatrix},
\qquad
g=Tq,
\]
其中 $S$ 与 $T$ 是第~\ref{sec:shrinking} 节中使用的共同关联映射与对比关联映射。令
\[
M^+=(M\ m),
\qquad
G^+=(TQ\ g),
\]
则对比维数 $\dim\Lambda$ 也恰好在下列条件下增加一：
\begin{equation}
m\in\operatorname{col}M
\quad\text{且}\quad
\begin{pmatrix}m\\g\end{pmatrix}
\notin
\operatorname{col}
\begin{pmatrix}M\\TQ\end{pmatrix}.
\label{eq:marginal-gap-condition}
\end{equation}
\end{theorem}

\begin{proof}
令
\[
A=
\begin{pmatrix}D\\Q\end{pmatrix},
\qquad
a=
\begin{pmatrix}d\\q\end{pmatrix}.
\]
新增一列会使 $\operatorname{rank}A$ 和 $\operatorname{rank}D$ 分别增加零或一。若 $d\notin\operatorname{col}D$，则 $a\notin\operatorname{col}A$，所以两个秩都增加一，其差不变。若 $d\in\operatorname{col}D$，则 $D$ 的秩不变，而 $A$ 的秩当且仅当 $a\notin\operatorname{col}A$ 时增加一。这就证明了式~\eqref{eq:marginal-info-01}--\eqref{eq:marginal-info-condition}。$A$--$B$ 的结论是把同一论证应用于 $(M,TQ)$。
\end{proof}

\begin{corollary}[比较数量的下界]
\label{cor:comparison-lower-bound}
若一个有限比较支持的当前投影有 $V$ 个当前节点、$E$ 条观测比较和 $c$ 个连通分量，则
\begin{equation}
\boxed{
\dim Q(\ker D)\le E-V+c.}
\label{eq:cycle-dimension-bound}
\end{equation}
因此，若某个支持在消去当前部分后携带 $k$ 个线性独立的延续方向，则必有
\begin{equation}
\boxed{
E\ge V-c+k.}
\label{eq:comparison-lower-bound}
\end{equation}
若当前图连通，则 $E\ge V-1+k$。
\end{corollary}

\begin{proof}
对节点--边关联矩阵，
\[
\operatorname{rank}D=V-c.
\]
因此
\[
\dim\ker D=E-V+c.
\]
由于 $Q(\ker D)$ 是 $\ker D$ 的像，其维数不可能超过式~\eqref{eq:cycle-dimension-bound}。整理即可得到式~\eqref{eq:comparison-lower-bound}。
\end{proof}

\begin{corollary}[单值延续支持]
\label{prop:single-valued-support}
假设在观测复合方案上，延续后果是当前节点的函数：存在映射 $\rho$，使得每个观测方案都具有形式 $(\xi,\rho(\xi))$。则
\[
Qr=0
\qquad
\text{对每个满足 }Dr=0\text{ 的 }r.
\]
因此，当前循环不能携带任何非零延续对比。
\end{corollary}

\begin{proof}
令 $R$ 把当前节点 $\xi$ 的基向量映射到延续标签 $\rho(\xi)$ 的基向量。对每条观测边，
\[
q_e=R(\mathbf e_{t(e)}-\mathbf e_{s(e)})=Rd_e.
\]
于是 $Q=RD$，而 $Dr=0$ 蕴含 $Qr=0$。
\end{proof}

因此，非零延续像要求命题~\ref{prop:support-rank} 中的秩确实增加。特别地，在相关观测支持上，延续后果不能是当前节点的单值函数。在贯穿全文的保险例子中，所需的非单点纤维只是：同一个当前合同后果允许出现不止一个续保条款。若未来后果由当前后果在技术上唯一决定，则该像为零。

令
\[
\gamma(z):=[\Gamma_A(z)]-[\Gamma_B(z)],
\qquad
\gamma(z)^{\mathsf T}v=H(z).
\]
为识别混合导数，设
\[
z_{00,n}=(x_0,a_0),\quad
z_{10,n}=(x_0+h_n,a_0),\quad
z_{01,n}=(x_0,a_0+k_n),\quad
z_{11,n}=(x_0+h_n,a_0+k_n),
\]
其中 $h_n,k_n\to0$ 且 $h_nk_n\ne0$。

选择四个当前后果
\[
\xi^R_{1n},\xi^R_{2n},\xi^R_{3n},\xi^R_{4n}
\]
并观察如下四个从左到右定向的补偿菜单：
\begin{align}
(\xi^R_{1n},\Gamma_B(z_{00,n}))
&\longrightarrow
(\xi^R_{2n},\Gamma_A(z_{00,n})),
\label{eq:primitive-rect-1}\\
(\xi^R_{2n},\Gamma_A(z_{10,n}))
&\longrightarrow
(\xi^R_{3n},\Gamma_B(z_{10,n})),
\label{eq:primitive-rect-2}\\
(\xi^R_{3n},\Gamma_B(z_{11,n}))
&\longrightarrow
(\xi^R_{4n},\Gamma_A(z_{11,n})),
\label{eq:primitive-rect-3}\\
(\xi^R_{4n},\Gamma_A(z_{01,n}))
&\longrightarrow
(\xi^R_{1n},\Gamma_B(z_{01,n})).
\label{eq:primitive-rect-4}
\end{align}
其延续对比分别为
\[
+\gamma(z_{00,n}),\quad
-\gamma(z_{10,n}),\quad
+\gamma(z_{11,n}),\quad
-\gamma(z_{01,n}).
\]

对行动方向的二阶导数，令
\[
z_{+,n}=(x_0+h_n,a_0),\qquad
z_{0}=(x_0,a_0),\qquad
z_{-,n}=(x_0-h_n,a_0),
\]
再选择四个当前后果
\[
\xi^X_{1n},\xi^X_{2n},\xi^X_{3n},\xi^X_{4n},
\]
并观察
\begin{align}
(\xi^X_{1n},\Gamma_B(z_{+,n}))
&\longrightarrow
(\xi^X_{2n},\Gamma_A(z_{+,n})),
\label{eq:primitive-curv-1}\\
(\xi^X_{2n},\Gamma_A(z_{0}))
&\longrightarrow
(\xi^X_{3n},\Gamma_B(z_{0})),
\label{eq:primitive-curv-2}\\
(\xi^X_{3n},\Gamma_A(z_{0}))
&\longrightarrow
(\xi^X_{4n},\Gamma_B(z_{0})),
\label{eq:primitive-curv-3}\\
(\xi^X_{4n},\Gamma_B(z_{-,n}))
&\longrightarrow
(\xi^X_{1n},\Gamma_A(z_{-,n})).
\label{eq:primitive-curv-4}
\end{align}
其延续对比分别为
\[
+\gamma(z_{+,n}),\quad
-\gamma(z_0),\quad
-\gamma(z_0),\quad
+\gamma(z_{-,n}).
\]

\begin{theorem}[双区块设计中的二阶识别]
\label{thm:derivatives-without-levels}
假设对每个 $n$，八个菜单
\eqref{eq:primitive-rect-1}--\eqref{eq:primitive-curv-4}
都被观察到；各设计中这些区块所使用的当前节点集合彼此不相交，并且没有其他观测比较使用这些当前节点。假设消去条件成立，并令 $H$ 在 $z_0$ 附近属于 $C^2$。

令 $b_n^R$ 与 $b_n^X$ 分别表示沿两个当前循环加总四个调整后精确转移所得的和。则
\begin{align}
b_n^R
&=
H(z_{00,n})-H(z_{10,n})-H(z_{01,n})+H(z_{11,n}),
\label{eq:rectangle-cycle-residual}\\
b_n^X
&=
H(z_{+,n})-2H(z_0)+H(z_{-,n}),
\label{eq:curvature-cycle-residual}
\end{align}
并且
\begin{equation}
\boxed{
\frac{b_n^R}{h_nk_n}\longrightarrow H_{xa}(z_0),
\qquad
\frac{b_n^X}{h_n^2}\longrightarrow H_{xx}(z_0).}
\label{eq:derivatives-without-levels-limit}
\end{equation}

任一抽样点上的水平值 $H(z)$ 都未被点识别。在约化式 $C^2$ 函数类上，观测为零的二阶 jet 恰好为
\begin{equation}
\boxed{
\left\{
(\theta_0,\theta_x,\theta_a,0,0,\theta_{aa})^{\mathsf T}
\right\}.}
\label{eq:block-null-jets}
\end{equation}
因此，一个局部二阶 jet 泛函得到点识别，当且仅当它是 $H_{xx}(z_0)$ 与 $H_{xa}(z_0)$ 的线性组合。
\end{theorem}

\begin{proof}
每组四个菜单的当前关联构成一个简单有向循环。把四条转移方程相加，会消去每一个当前节点价值。式~\eqref{eq:primitive-rect-1}--\eqref{eq:primitive-rect-4} 中的延续对比加总为式~\eqref{eq:rectangle-cycle-residual}；式~\eqref{eq:primitive-curv-1}--\eqref{eq:primitive-curv-4} 中的延续对比加总为式~\eqref{eq:curvature-cycle-residual}。式~\eqref{eq:derivatives-without-levels-limit} 中的极限就是通常的混合矩形差分与中心二阶差分。

令
\[
q(x,a)=\alpha+\beta(x-x_0)+\psi(a),
\]
其中 $\alpha,\beta\in\mathbb R$ 且 $\psi\in C^2$。它的矩形差分与关于 $x$ 的中心二阶差分都为零。因此，把 $H$ 替换成 $H+q$ 并不会改变任何观测当前循环上的延续和。由于每个当前区块都是简单循环，且不同区块使用互不相交的当前节点，可以沿前三条边依次调整当前节点价值；第四条边之所以也成立，恰好是因为 $q$ 的循环和为零。因此 $H$ 与 $H+q$ 能够复现每一个原始转移观测。特别地，没有任何抽样水平 $H(z)$ 得到点识别。

反过来，每个观测为零的扰动对每个 $n$ 都有标准化矩形残差与曲率残差为零，因此
\[
q_{xa}(z_0)=q_{xx}(z_0)=0.
\]
任意这两个分量为零的二阶 jet，都可以由
\[
q(x,a)
=
\alpha+\beta(x-x_0)+\gamma(a-a_0)
+\frac{\delta}{2}(a-a_0)^2
\]
实现，而根据前述论证，这样的 $q$ 在观测上为零。这就证明了式~\eqref{eq:block-null-jets}。两个有限差分极限识别 $xx$ 与 $xa$ 方向，而零扰动则表明不存在其他被识别的二阶 jet 方向。
\end{proof}

\begin{corollary}[双区块设计中的边最小性]
\label{cor:two-block-edge-minimal}
若四个矩形菜单位置中的任意一个在所有设计中都缺失，则无法由剩余双区块数据点识别 $H_{xa}(z_0)$。若四个曲率菜单位置中的任意一个在所有设计中都缺失，则无法点识别 $H_{xx}(z_0)$。
\end{corollary}

\begin{proof}
从每个矩形区块删除任意一个菜单后，其四个当前节点之间只剩一棵树，因此该区块不存在非零循环。剩余曲率区块无法识别 $H_{xa}$：扰动
\[
q(x,a)=(x-x_0)(a-a_0)
\]
的曲率残差为零，却改变 $q_{xa}(z_0)$。由于不完整的矩形区块是树，可以通过调整当前节点价值保持其边方程不变，因此 $H_{xa}(z_0)$ 未识别。类似地，若从每个曲率区块删除一个菜单，矩形区块看不到
\[
q(x,a)=(x-x_0)^2,
\]
而不完整曲率区块又都是树，因此 $H_{xx}(z_0)$ 未识别。
\end{proof}

每一个导数都由一个四边当前循环承载，而延续水平值本身仍不受限制。

\begin{proposition}[额外比较]
\label{prop:additional-comparisons}
假设一个扩展的局部比较支持包含定理~\ref{thm:derivatives-without-levels} 中的两个四边区块。令 $\widetilde\Lambda_n$ 表示扩展设计中所有观测当前循环所生成的 $A$--$B$ 对比空间。假设每个 $\lambda\in\widetilde\Lambda_n$ 都满足
\begin{align}
\sum_{i:a_{in}=\bar a}\lambda_i&=0
\quad\text{对每个抽样状态 }\bar a,
\label{eq:additional-state-balance}\\
\sum_i\lambda_i(x_{in}-x_0)&=0.
\label{eq:additional-x-balance}
\end{align}
则在约化式 $C^2$ 函数类上，被识别的二阶 jet 泛函仍然恰好是
\begin{equation}
\boxed{\operatorname{span}\{\partial_{xx},\partial_{xa}\}.}
\label{eq:robust-identified-space}
\end{equation}
特别地，与响应有关的两个导数仍得到点识别，而没有任何水平值 $H(z)$ 得到点识别。额外比较可以与原区块共享当前节点。
\end{proposition}

\begin{proof}
原来的两个区块仍被观测，因此每个观测等价扰动都满足
\[
q_{xx}(z_0)=q_{xa}(z_0)=0.
\]
反过来，对
\[
q(x,a)=\alpha+\beta(x-x_0)+\psi(a),
\qquad \psi\in C^2,
\]
条件~\eqref{eq:additional-state-balance}--\eqref{eq:additional-x-balance} 给出
\[
\sum_i\lambda_iq(z_{in})=0
\qquad
\forall\lambda\in\widetilde\Lambda_n.
\]
所以这些扰动在扩展循环系统中仍然观测为零。其二阶 jet 恰好穷尽所有 $xx$ 与 $xa$ 分量为零的向量。结合原有矩形差分和中心二阶差分，就得到式~\eqref{eq:robust-identified-space}。常数扰动则表明不存在被识别的点水平值。
\end{proof}

\subsection{高阶局部信息}
\label{sec:higher-order}

上述论证并不限于二阶 jet。令 $m\ge1$ 且
\[
\alpha=(\alpha_1,\alpha_2)\in\mathbb N_0^2,
\qquad
|\alpha|=\alpha_1+\alpha_2,
\qquad
\alpha!=\alpha_1!\alpha_2!.
\]
对
\[
\delta z_{in}
=
(\delta x_{in},\delta a_{in})
=
z_{in}-z_0,
\]
写
\[
(\delta z_{in})^\alpha
=
\delta x_{in}^{\alpha_1}
\delta a_{in}^{\alpha_2}.
\]
令 $t=(t_\alpha)_{|\alpha|\le m}$，并定义
\begin{equation}
\mathcal L_t^{(m)}f(z_0)
:=
\sum_{|\alpha|\le m}
t_\alpha\,\partial^\alpha f(z_0).
\label{eq:higher-order-functional}
\end{equation}

\begin{theorem}[高阶有限差分]
\label{thm:higher-order-differences}
假设局部设计的支持半径趋于零。令 $\lambda_n\in\Lambda_n$ 且 $s_n\ne0$，并对每个多重指标 $|\alpha|\le m$ 满足
\begin{equation}
\frac1{s_n}
\sum_i
\lambda_{in}
\frac{(\delta z_{in})^\alpha}{\alpha!}
\longrightarrow
t_\alpha,
\label{eq:higher-order-moments}
\end{equation}
以及
\begin{equation}
\frac1{|s_n|}
\sum_i
|\lambda_{in}|
\|\delta z_{in}\|^m
=
O(1).
\label{eq:higher-order-variation}
\end{equation}
则对 $z_0$ 附近任意 $f\in C^m$，
\begin{equation}
\boxed{
\frac1{s_n}\sum_i\lambda_{in}f(z_{in})
\longrightarrow
\mathcal L_t^{(m)}f(z_0).}
\label{eq:higher-order-limit}
\end{equation}
在消去条件下，同一极限可由相应的标准化循环和识别。
\end{theorem}

\begin{proof}
在 $z_0$ 处使用 Taylor 公式，并且在收缩支持上一致地有
\[
f(z_{in})
=
\sum_{|\alpha|\le m}
\frac{\partial^\alpha f(z_0)}{\alpha!}
(\delta z_{in})^\alpha
+
R_{in},
\qquad
R_{in}=o(\|\delta z_{in}\|^m).
\]
乘以 $\lambda_{in}/s_n$ 并求和后，多项式部分由式~\eqref{eq:higher-order-moments} 收敛；余项由式~\eqref{eq:higher-order-variation} 趋于零。当 $f=H$ 时，左端由式~\eqref{eq:gap-cycle-price} 等于标准化的观测循环转移。
\end{proof}

令 $\Pi_{m-1}$ 表示 $(x,a)$ 中总次数不超过 $m-1$ 的多项式集合。

\begin{theorem}[不恢复低阶信息的高阶识别]
\label{thm:order-separation}
假设每一个观测 $A$--$B$ 循环对比都会消去低阶多项式：
\begin{equation}
\sum_i\lambda_i p(z_{in})=0
\qquad
\text{对每个 }n,\ \lambda\in\Lambda_n,\ p\in\Pi_{m-1}.
\label{eq:lower-order-annihilation}
\end{equation}
再假设定理~\ref{thm:higher-order-differences} 对某个系数向量 $t$ 成立，并且
\[
t_\alpha=0
\qquad\text{只要 }|\alpha|<m.
\]
那么齐次 $m$ 阶泛函
\begin{equation}
\boxed{
\sum_{|\alpha|=m}
t_\alpha\,\partial^\alpha H(z_0)}
\label{eq:mth-order-identified}
\end{equation}
得到点识别；但在约化式 $C^m$ 函数类上，任何只涉及严格低于 $m$ 阶导数的非零局部泛函都未被点识别。

更一般地，若这类有限差分权重序列张成齐次 $m$ 阶导数空间中的一个子空间 $\mathcal S_m$，则 $\mathcal S_m$ 中的每一个元素都得到识别，即使完整的 $(m-1)$ 阶 jet 仍未识别。
\end{theorem}

\begin{proof}
式~\eqref{eq:mth-order-identified} 的识别来自定理~\ref{thm:higher-order-differences}。由式~\eqref{eq:lower-order-annihilation}，每个 $p\in\Pi_{m-1}$ 都在观测上为零。对 $(m-1)$ 阶 jet 的任意非零线性泛函，可取其 Taylor 多项式 $p\in\Pi_{m-1}$，使该泛函在 $z_0$ 处非零。加入任意倍数的 $p$ 都不会改变任何循环方程，却能让目标值遍历整个实数轴。因此，不存在被点识别的非零低阶 jet 泛函。同一论证可同时应用于所有张成 $\mathcal S_m$ 的权重序列。
\end{proof}

定理~\ref{thm:derivatives-without-levels} 中的二阶设计就是 $m=2$ 的情形。其循环权重消去仿射函数，而其二阶矩张成
\[
\operatorname{span}\{\partial_{xx},\partial_{xa}\}.
\]
额外的零方向 $\partial_{aa}$ 由该特定设计的支持几何所决定。

\begin{corollary}[不恢复水平值的反事实比较静态]
\label{cor:response-without-levels}
在定理~\ref{thm:derivatives-without-levels} 的条件下，假设环境 $A$ 中的补偿观测识别 $W_{A,xa}$ 与 $W_{A,xx}$，已知当前侧导数 $d_{xa}$ 与 $d_{xx}$，并且目标基准是一个内部严格局部最优点。则 $D_ag_B(a_0)$ 得到点识别，即使任何邻近设计点上的 $H(z)$ 都未得到点识别。
\end{corollary}

\begin{proof}
该定理识别 $\eta_{xa}=H_{xa}(z_0)$ 与 $\eta_{xx}=H_{xx}(z_0)$。把它们代入传递恒等式与目标一阶条件，即得式~\eqref{eq:main-response-v19}。
\end{proof}

\section{反事实比较静态}
\label{sec:counterfactual}

令 $A$ 为已观测的锚环境，$B$ 为目标环境。写为
\begin{equation}
W_J(x,a)=u_J(x,a)+V(\Gamma_J(x,a)),
\qquad J\in\{A,B\},
\label{eq:WJ-decomp-v19}
\end{equation}
其中当前侧差异 $u_B-u_A$ 已知。令
\begin{align}
d_{xa}&:=(u_B-u_A)_{xa},
\label{eq:d-xa-v19}\\
d_{xx}&:=(u_B-u_A)_{xx}.
\label{eq:d-xx-v19}
\end{align}
由于 $H=V\circ\Gamma_A-V\circ\Gamma_B$，第~\ref{sec:shrinking} 节识别的修正项满足
\begin{align}
\boxed{W_{B,xa}=W_{A,xa}+d_{xa}-\eta_{xa},}
\label{eq:transport-xa-v19}\\
\boxed{W_{B,xx}=W_{A,xx}+d_{xx}-\eta_{xx}.}
\label{eq:transport-xx-v19}
\end{align}

令 $C_B(x,a)$ 为目标环境中的货币资源成本。观测到目标基准选择 $x=g_B(a_0)$。假设它是内部严格局部最优点，并且在所考察点上
\begin{equation}
C_{B,xx}-W_{B,xx}>0.
\label{eq:target-regularity-v19}
\end{equation}

\begin{theorem}[反事实比较静态]
\label{thm:counterfactual-v19}
在 $(g_B(a_0),a_0)$ 处，假设环境 $A$ 中的补偿观测识别 $W_{A,xa}$ 与 $W_{A,xx}$，当前侧导数 $d_{xa}$ 与 $d_{xx}$ 已知，并且局部比较设计包含满足式~\eqref{eq:moment-limit}--\eqref{eq:moment-variation}、分别对应 $\partial_{xa}$ 与 $\partial_{xx}$ 的有限差分权重。则 $D_ag_B(a_0)$ 得到点识别，且
\begin{equation}
\boxed{
D_ag_B(a_0)
=
\frac{
W_{A,xa}+d_{xa}-\eta_{xa}-C_{B,xa}
}{
C_{B,xx}-W_{A,xx}-d_{xx}+\eta_{xx}
}.}
\label{eq:main-response-v19}
\end{equation}
\end{theorem}

\begin{proof}
目标环境的一阶条件为
\[
W_{B,x}(g_B(a),a)=C_{B,x}(g_B(a),a).
\]
沿正则局部分支求导，得到
\[
(C_{B,xx}-W_{B,xx})D_ag_B
=
W_{B,xa}-C_{B,xa}.
\]
利用式~\eqref{eq:transport-xa-v19}--\eqref{eq:transport-xx-v19} 以及定理~\ref{thm:primitive-local} 识别 $\eta_{xa}$ 与 $\eta_{xx}$。
\end{proof}

当环境 $A$ 中可以观察邻近补偿比较时，式~\eqref{eq:pi-derivatives} 给出
\begin{equation}
W_{A,xa}=\pi_{A,a},
\qquad
W_{A,xx}=\pi_{A,x},
\label{eq:anchor-pi-derivs-v19}
\end{equation}
于是
\begin{equation}
\boxed{
D_ag_B(a_0)
=
\frac{
\pi_{A,a}+d_{xa}-\eta_{xa}-C_{B,xa}
}{
C_{B,xx}-\pi_{A,x}-d_{xx}+\eta_{xx}
}.}
\label{eq:anchor-to-target-v19}
\end{equation}
当环境 $B$ 中也可以观察邻近补偿比较时，式~\eqref{eq:direct-gprime} 可直接使用。

\subsection{高阶比较静态}
\label{sec:policy-hierarchy}

令
\[
F_B(x,a)=W_{B,x}(x,a)-C_{B,x}(x,a),
\qquad
F_B(g_B(a),a)=0.
\]
重复使用隐式求导是标准做法。一旦 $D_ag_B(a_0),\ldots,D_a^{r-1}g_B(a_0)$ 已知，进入 $D_a^rg_B(a_0)$ 的唯一新的 $(r+1)$ 阶延续项是方向泛函
\begin{equation}
\boxed{
\left(\partial_a+D_ag_B(a_0)\partial_x\right)^r H_x(z_0),
\qquad z_0=(g_B(a_0),a_0),}
\label{eq:response-correction-directional}
\end{equation}
其中微分算子中的系数 $D_ag_B(a_0)$ 保持固定。因此，只要循环矩生成式~\eqref{eq:response-correction-directional}，并结合已知的参考环境导数、当前侧导数和成本导数，定理~\ref{thm:higher-order-differences} 就识别第 $r$ 阶政策导数。无需恢复完整的 $(r+1)$ 阶延续 jet。例如，为识别 $D_a^2g_B(a_0)$ 所需的新三阶延续项为
\[
H_{xaa}+2(D_ag_B)H_{xxa}+(D_ag_B)^2H_{xxx}.
\]
因此，第~\ref{sec:higher-order} 节中的阶次分离直接传递到高阶比较静态。

\section{多个环境的扩展}
\label{sec:environment-network}

令 $\mathcal J$ 为有限环境集合，并写
\[
\Psi_J=V\circ\Gamma_J.
\]
对任何已经由比较数据识别的线性局部泛函 $\mathcal L$，定义观测到的成对修正
\[
\eta_{JK}^{\mathcal L}=\mathcal L(\Psi_J-\Psi_K).
\]
这些量就是以各环境为节点的图上的普通势差。

\begin{proposition}[路径无关性]
\label{thm:path-independence}
令 $\mathsf B$ 为环境图的节点--边关联矩阵，并令 $\eta^{\mathcal L}$ 收集带方向的边修正。下列条件等价：
\begin{enumerate}[label=(\roman*)]
\item 存在节点势 $\theta_J$，使每条观测边上都有 $\eta_{JK}^{\mathcal L}=\theta_J-\theta_K$；
\item 对每个 $c\in\ker\mathsf B$，都有 $c^{\mathsf T}\eta^{\mathcal L}=0$；
\item 两个连通环境之间的路径和与所选路径无关。
\end{enumerate}
在共同延续表示下，可以取 $\theta_J=\mathcal L\Psi_J$。因此，一棵生成树就足以恢复一个连通分量中任意参考环境到目标环境的修正，而每一条独立的非树边都会提供一个循环型过度识别限制。
\end{proposition}

\begin{proof}
第一个条件等价于 $\eta^{\mathcal L}=\mathsf B^{\mathsf T}\theta$。由于
$\operatorname{col}(\mathsf B^{\mathsf T})=(\ker\mathsf B)^\perp$，它与第二个条件等价。路径结论由望远镜相消立即得到。
\end{proof}

这个扩展并没有引入第二种识别机制。当前比较图识别边修正；环境图只是在此基础上组织已经识别的势差。若某个目标环境位于参考环境连通分量之外，则仅凭这些边差异无法确定从参考环境到该目标环境的修正。

\section{一个动态后果}
\label{sec:dynamics}

令继承状态为 $s\in\R^d$，当前行动为 $x\in\R^m$，已知的技术运动规律为
\[
s^+=T(s,x).
\]
若比较分析识别一个可微政策 $x=g(s)$，则诱导出的闭环映射为
\begin{equation}
\Phi(s):=T(s,g(s)).
\label{eq:closed-loop-map}
\end{equation}
由链式法则，
\begin{equation}
\boxed{
D\Phi(s)=D_sT(s,g(s))+D_xT(s,g(s))Dg(s).}
\label{eq:closed-loop-jacobian}
\end{equation}
因此，识别政策导数就识别了状态转移的 Jacobian。只要相应的政策 jet 已识别，更高阶的转移导数由普通多元链式法则得到。

从这里开始，结论都是标准的离散时间动力系统结果 \citep{elaydi2005,kuznetsov2004}。在双曲不动点处，式~\eqref{eq:closed-loop-jacobian} 的乘子给出通常的局部稳定性和鞍点结论。在临界乘子处，更高阶转移导数可能变得重要。对于归一化到零且满足 $\Phi'(0)=-1$ 的标量不动点，通常的三次 flip 系数为
\begin{equation}
\boxed{
c_{\mathrm{flip}}
=
\frac14\Phi''(0)^2+\frac16\Phi'''(0).}
\label{eq:flip-coefficient}
\end{equation}
等价地，
\[
\Phi^2(y)=y-2c_{\mathrm{flip}}y^3+o(|y|^3).
\]
因此，只要比较数据识别三阶转移 jet，它们也就识别这个标准系数。一对非实的单位模乘子构成标准的 Neimark--Sacker 临界配置；同样地，本文真正关心的识别问题是转移导数与正规形系数，而不是动力学分类本身。图~\ref{fig:critical-dynamics} 概括了这一区别。

\begin{figure}[H]
\centering
\refstepcounter{figure}\label{fig:critical-dynamics}
\textbf{图 \thefigure. Jacobian 之外的临界动力学}\par\smallskip
\begin{minipage}[c]{0.755\textwidth}
\centering
\includegraphics[width=\linewidth,trim=0 0 410bp 78bp,clip]{fig4_critical.jpg}
\end{minipage}\hfill
\begin{minipage}[c]{0.225\textwidth}
\centering\small
\textbf{识别阶梯}\par\medskip
\fbox{\parbox{0.84\linewidth}{\centering 补偿比较}}\\[4pt]
$\downarrow$\\[4pt]
\fbox{\parbox{0.84\linewidth}{\centering 延续导数}}\\[4pt]
$\downarrow$\\[4pt]
\fbox{\parbox{0.84\linewidth}{\centering 政策 jet}}\\[4pt]
$\downarrow$\\[4pt]
\fbox{\parbox{0.84\linewidth}{\centering 转移 jet}}\\[7pt]
\fcolorbox{blue!60!black}{white}{\parbox{0.86\linewidth}{\centering\scriptsize
无单位模乘子：\\一阶线性化即可分类}}\\[7pt]
\fcolorbox{orange!80!black}{white}{\parbox{0.86\linewidth}{\centering\scriptsize
存在单位模乘子：\\高阶漂移可能决定结果}}
\end{minipage}
\par\smallskip\small\textit{图示说明：左侧原图展示一维 flip 与二维临界旋转中的稳定性边界及高阶漂移；右侧强调本文的识别顺序：由补偿比较得到延续导数，再得到政策 jet 与转移 jet。远离临界边界时一阶 Jacobian 足够；处于临界区间时高阶项可能决定局部动力学。}
\end{figure}

在平稳贴现的特例下，可以把标准 Bellman 验证与价值迭代应用于已经恢复的政策 \citep{bellman1957,stokeylucasprescott1989}。因此，递归在识别之后只是表示与计算工具，而不是识别限制。

\section{适用域与解释}

货币原语是一个可转移的计价物：精确补偿转移可以在同一个加法尺度上比较。消去条件决定观测转移是否具有加法表示~\eqref{eq:graph-representation}。其中的势是有限数据的表示坐标，而不是第~\ref{sec:locus} 节结构分解中预先假定的价值。

目标环境与观测比较族是两个不同对象。即使数据中缺少某些补偿比较，目标延续映射仍在其局部选择集上有定义。贯穿局部结果的“识别”若没有额外结构限定，都是指满足观测循环方程的约化式光滑延续函数类。固定共同 $V$ 以及固定延续映射 $(\Gamma_A,\Gamma_B)$ 所生成的限制，可能进一步缩小结构识别集。

两阶段保险例子隔离出真正重要的支持条件：可以保持当前保险后果不变，同时在诱导菜单中改变状态依赖续保条款，并通过保费转移重新建立无差异。这样的合同变化会产生一个非单点延续纤维。反之，若未来后果在技术上由当前后果唯一决定，推论~\ref{prop:single-valued-support} 给出零延续像。因此，这个结果区分的是可分配的延续合同与技术上固定的延续，而不是把“重组”当作所有动态选择问题都具备的性质。

光滑结果使用收缩循环关系、Taylor 展开与极限。高阶比较静态来自对一阶条件的重复求导。环境图若被使用，只包含当前比较图已经识别出来的势差。动态部分则把已经识别的政策与技术运动规律复合；后续所有稳定性与递归结论，都是所恢复政策与转移导数的标准后果。

\section{结论}

一次精确补偿比较把有限的当前变化置于共同货币尺度上。当目标环境本身在局部被观察时，在邻近状态上重复实验即可得到直接比较静态。若目标环境在局部没有这些比较，则观测比较图中的循环会消去不受限制的当前节点价值，只留下延续差异。

一个比较设计的有限信息内容是延续像 $Q(\ker D)$。其维数是一个秩差；只有当一条新比较在当前侧冗余、而在当前--延续联合侧不冗余时，它才增加一个信息方向。这把“当前投影闭合而延续提升不闭合”的图形直觉转化成逐边的信息陈述。

在局部，收缩循环权重能够在不识别水平值的情况下识别导数。在双区块设计中，$H_{xa}$ 与 $H_{xx}$ 得到点识别，而所有抽样值 $H(z)$ 仍未识别；从相应区块删除一条边，就会破坏对应导数的识别。相同论证可推广到任意导数阶次：低阶多项式项可以在观测上为零，而某个高阶矩仍然存活。

这些已识别延续泛函进入目标一阶条件，并决定反事实比较静态。支持条件具有直接的合同解释：同一个当前后果必须有时允许对应不止一个被观察到的延续条款。随后，可以在这些已恢复对象之上使用跨环境势差和标准状态空间动力学，而不需要再增加新的识别限制。

至此，这一叙述可以恰好闭合；余下的问题，是之后所发生之事的代价。

\appendix

\section{有限图定理的证明}
\label{app:finite-proof}

\begin{lemma}[有理零空间基]
\label{lem:rational-v19}
若一个矩阵的元素都是有理数，则它的零空间存在一组坐标为有理数的基。
\end{lemma}

\begin{proof}
在 $\mathbb Q$ 上做行消元，可把主元变量表示为自由变量的有理线性组合。依次把自由变量设为标准基向量，即得到有理基。
\end{proof}

\begin{lemma}[清除分母]
\label{lem:clearing-v19}
令 $M$ 为整数矩阵，$\ell$ 为整数向量。若 $\ell$ 属于 $M$ 的实列空间，则存在 $k\in\N_+$ 与 $z\in\Z^{\dim(\text{domain }M)}$，使得
\[
k\ell=Mz.
\]
\end{lemma}

\begin{proof}
方程 $Mz=\ell$ 存在实数解。在 $\mathbb Q$ 上做 Gaussian 消元可得到有理解；再乘以共同分母即可。
\end{proof}

\begin{proof}[定理~\ref{thm:additive-graph} 的证明]
令
\[
M=
\begin{pmatrix}
D\\Q
\end{pmatrix}.
\]
若式~\eqref{eq:graph-representation} 成立且 $Mz=0$，则 $z^{\mathsf T}\tau=0$，故消去条件是必要的。

反过来，$M$ 的元素均为整数。由引理~\ref{lem:rational-v19}，$\ker M$ 中的每个实向量都是有理零空间向量的实线性组合；清除分母后，这些向量都可以化为整数零向量。因此消去条件意味着
\[
Mz=0\Longrightarrow z^{\mathsf T}\tau=0
\qquad(z\in\R^E).
\]
于是 $\tau\perp\ker M$，所以 $\tau$ 属于 $M$ 的行空间，等价地属于 $M^{\mathsf T}$ 的列空间。因此存在 $u$ 与 $v$，使
\[
\tau=M^{\mathsf T}(u,v)=D^{\mathsf T}u+Q^{\mathsf T}v.
\]
\end{proof}

\begin{proof}[定理~\ref{thm:feasible-finite} 的证明]
令 $w=(u,v)$，并令 $M=(D^{\mathsf T},Q^{\mathsf T})$，于是观测方程为 $Mw=\tau$。目标泛函为
\[
\ell^{\mathsf T}w,
\qquad
\ell=(0,c).
\]
该泛函在所有 $Mw=\tau$ 的解上保持常数，当且仅当 $\ell$ 与 $\ker M$ 正交，等价地，当且仅当 $\ell$ 属于 $M$ 的行空间。因此存在实边向量 $z$，使
\[
Dz=0,
\qquad
Qz=c.
\]
由于 $D,Q,c$ 均为整数，引理~\ref{lem:clearing-v19} 给出有限复制形式~\eqref{eq:feasible-cycle}。用 $z^{\mathsf T}$ 左乘式~\eqref{eq:graph-representation} 即得式~\eqref{eq:feasible-price}。

若 $\ell$ 不属于 $M$ 的行空间，则存在扰动 $h\in\ker M$ 使 $\ell^{\mathsf T}h\ne0$。从任意一个解出发，$w+th$ 对所有 $t\in\R$ 都保持观测等价，而目标泛函则遍历整个 $\R$。
\end{proof}

\begin{thebibliography}{99}

\bibitem[Afriat(1967)]{afriat1967}
Afriat, S. N. (1967).
The construction of utility functions from expenditure data.
\emph{International Economic Review} 8, 67--77.

\bibitem[Bellman(1957)]{bellman1957}
Bellman, R. (1957).
\emph{Dynamic Programming}.
Princeton, NJ: Princeton University Press.

\bibitem[Elaydi(2005)]{elaydi2005}
Elaydi, S. (2005).
\emph{An Introduction to Difference Equations}, 3rd ed.
New York: Springer.

\bibitem[Bhattacharya(2015)]{bhattacharya2015}
Bhattacharya, D. (2015).
Nonparametric welfare analysis for discrete choice.
\emph{Econometrica} 83, 617--649.

\bibitem[Brown and Calsamiglia(2007)]{browncalsamiglia2007}
Brown, D. J. and C. Calsamiglia (2007).
The nonparametric approach to applied welfare analysis.
\emph{Economic Theory} 31, 183--188.

\bibitem[Castillo and Freer(2023)]{castillofreer2023}
Castillo, M. and M. Freer (2023).
A general revealed preference test for quasilinear preferences: theory and experiments.
\emph{Experimental Economics} 26, 673--696.

\bibitem[Chambers, Echenique, and Lambert(2021)]{chambersetal2021}
Chambers, C. P., F. Echenique, and N. S. Lambert (2021).
Recovering preferences from finite data.
\emph{Econometrica} 89, 1633--1664.

\bibitem[Debreu(1954)]{debreu1954}
Debreu, G. (1954).
Representation of a preference ordering by a numerical function.
In R. M. Thrall, C. H. Coombs, and R. L. Davis (eds.), \emph{Decision Processes}, 159--165. New York: Wiley.

\bibitem[Doignon and Falmagne(1974)]{doignonfalmagne1974}
Doignon, J. P. and J. C. Falmagne (1974).
Difference measurement and simple scalability with restricted solvability.
\emph{Journal of Mathematical Psychology} 11, 473--499.

\bibitem[Fishburn(2001)]{fishburn2001}
Fishburn, P. C. (2001).
Cancellation conditions for finite two-dimensional additive measurement.
\emph{Journal of Mathematical Psychology} 45, 2--26.

\bibitem[Fishburn, Marcus-Roberts, and Roberts(1988)]{fishburnetal1988}
Fishburn, P. C., H. M. Marcus-Roberts, and F. S. Roberts (1988).
Unique finite difference measurement.
\emph{SIAM Journal on Discrete Mathematics} 1, 334--354.

\bibitem[Fuchs-Seliger(1989)]{fuchsseliger1989}
Fuchs-Seliger, S. (1989).
Money-metric utility functions in the theory of revealed preference.
\emph{Mathematical Social Sciences} 18, 199--210.

\bibitem[Fuchs-Seliger(1990)]{fuchsseliger1990}
Fuchs-Seliger, S. (1990).
An axiomatic approach to compensated demand.
\emph{Journal of Economic Theory} 52, 111--122.

\bibitem[Gorno(2019)]{gorno2019}
Gorno, L. (2019).
Revealed preference and identification.
\emph{Journal of Economic Theory} 183, 698--739.

\bibitem[Kraft, Pratt, and Seidenberg(1959)]{kraftprattseidenberg1959}
Kraft, C. H., J. W. Pratt, and A. Seidenberg (1959).
Intuitive probability on finite sets.
\emph{Annals of Mathematical Statistics} 30, 408--419.

\bibitem[Kuznetsov(2004)]{kuznetsov2004}
Kuznetsov, Y. A. (2004).
\emph{Elements of Applied Bifurcation Theory}, 3rd ed.
New York: Springer.

\bibitem[Makowski and Ostroy(2013)]{makowskiostroy2013}
Makowski, L. and J. M. Ostroy (2013).
From revealed preference to preference revelation.
\emph{Journal of Mathematical Economics} 49, 71--81.

\bibitem[Nocke and Schutz(2017)]{nockeschutz2017}
Nocke, V. and N. Schutz (2017).
Quasi-linear integrability.
\emph{Journal of Economic Theory} 169, 603--628.

\bibitem[Rochet(1987)]{rochet1987}
Rochet, J.-C. (1987).
A necessary and sufficient condition for rationalizability in a quasi-linear context.
\emph{Journal of Mathematical Economics} 16, 191--200.

\bibitem[Samuelson(1938)]{samuelson1938}
Samuelson, P. A. (1938).
A note on the pure theory of consumer's behaviour.
\emph{Economica} 5, 61--71.

\bibitem[Samuelson(1948)]{samuelson1948}
Samuelson, P. A. (1948).
Consumption theory in terms of revealed preference.
\emph{Economica} 15, 243--253.

\bibitem[Samuelson(1950)]{samuelson1950}
Samuelson, P. A. (1950).
The problem of integrability in utility theory.
\emph{Economica} 17, 355--385.

\bibitem[Scott(1964)]{scott1964}
Scott, D. (1964).
Measurement structures and linear inequalities.
\emph{Journal of Mathematical Psychology} 1, 233--247.

\bibitem[Stokey, Lucas, and Prescott(1989)]{stokeylucasprescott1989}
Stokey, N. L., R. E. Lucas, Jr., and E. C. Prescott (1989).
\emph{Recursive Methods in Economic Dynamics}.
Cambridge, MA: Harvard University Press.

\end{thebibliography}

\end{document}

